\section{Step 67 Repeat a previous length only when one exists}\label{png-huff-s67}
Implement \texttt{huff\_repeat16}. If the decoded operation is not sixteen, select the zero-repeat dispatch. For sixteen, reject a cursor still at the start of the combined vector. Load the previous expanded length, retain it across a two-bit read, and prepare a run of field plus three entries.

Use the repeat16 fixture. A direct length one and then a direct zero precede this operation. Extra field three requests six more zeros. Repeating zero is permitted because a previous length exists; the error is using operation sixteen before any previous entry, not repeating an absent symbol's zero length.

\paragraph{Observe} Use \texttt{s67-1}; stop before \texttt{huff\_repeat16\_ready}. Inspect the state described above before continuing.

\Needspace{7\baselineskip}
\paragraph{Check} Can operation sixteen repeat a previous zero length?
\begin{itemize}
  \item[A.] Yes, when a previous expanded entry exists.
  \item[B.] No, it can repeat only nonzero lengths.
  \item[C.] Only in a fixed block.
\end{itemize}
